
\documentclass{article}
\usepackage{colortbl}
\usepackage{makecell}
\usepackage{multirow}
\usepackage{supertabular}

\begin{document}

\newcounter{utterance}

\centering \large Interaction Transcript for game `ifeval', experiment `ifeval', episode 367 with llama{-}3.1{-}8b{-}instruct.
\vspace{24pt}

{ \footnotesize  \setcounter{utterance}{1}
\setlength{\tabcolsep}{0pt}
\begin{supertabular}{c@{$\;$}|p{.15\linewidth}@{}p{.15\linewidth}p{.15\linewidth}p{.15\linewidth}p{.15\linewidth}p{.15\linewidth}}
    \# & \multicolumn{2}{c}{Player} && \multicolumn{2}{c}{Game Master} \\
    \hline

    \theutterance \stepcounter{utterance}  
    & & & \multicolumn{4}{p{0.6\linewidth}}{
        \cellcolor[rgb]{0.9,0.9,0.9}{
            \makecell[{{p{\linewidth}}}]{
                \texttt{\tiny{[P1$\langle$GM]}}
                \texttt{Expand the riddle into a story with a funny tone:} \\
\\ 
\texttt{What can you catch but not throw?} \\
\texttt{A cold} \\
\\ 
\texttt{Use * to highlight at least 2 sections in your text. For example: *this is a highlighted text section*.} \\
            }
        }
    }
     \\ \\

    \theutterance \stepcounter{utterance}  
    & \multicolumn{4}{p{0.6\linewidth}}{
        \cellcolor[rgb]{0.561,0.737,0.561}{
            \makecell[{{p{\linewidth}}}]{
                \texttt{\tiny{[P1$\rangle$GM]}}
                \texttt{Once upon a time, in the quirky town of Sneezeville, there lived a lovable but slightly clumsy fellow named Bob. Bob was known for his impressive collection of colorful scarves and his ability to trip over air molecules. One day, while out for a stroll, Bob stumbled upon a mysterious sign that read: "What can you catch but not throw?"} \\
\\ 
\texttt{Intrigued, Bob decided to embark on a quest to find the answer. He asked his best friend, a wise{-}cracking chicken named Cluck Norris (yes, that's his real name), for help. Cluck thought for a moment and then said, *"Dude, it's a cold! You can catch a cold, but you can't throw one."*} \\
\\ 
\texttt{Bob was taken aback by Cluck's clever response. "Who knew chickens were so smart?" he exclaimed. But then he started to think about it... and realized that Cluck was absolutely right. You can catch a cold by being near someone who has one, but you can't throw a cold at someone like a frisbee.} \\
\\ 
\texttt{As Bob continued his journey, he encountered all sorts of wacky characters who tried to stump him with their own versions of the riddle. There was the sneezy librarian who said, "It's a tissue!" (but Bob knew you could throw a tissue). There was the mischievous magician who said, "It's a magic trick!" (but Bob wasn't buying it). And there was even the grumpy old man who said, "It's a cold, but only if you're a grumpy old man!" (which, let's be honest, is a pretty good answer).} \\
\\ 
\texttt{In the end, Bob returned to Sneezeville a hero, with Cluck Norris by his side. Together, they celebrated their victory with a big bowl of chicken soup (which, of course, was made with love and a dash of Cluck's secret spice). And from that day on, Bob and Cluck were known as the dynamic duo of riddle{-}solving, always ready to catch a cold (just kidding, they were just ready to solve more riddles).} \\
            }
        }
    }
    & & \\ \\

    \theutterance \stepcounter{utterance}  
    & & & \multicolumn{2}{p{0.3\linewidth}}{
        \cellcolor[rgb]{0.9,0.9,0.9}{
            \makecell[{{p{\linewidth}}}]{
                \texttt{\tiny{[GM$|$GM]}}
                {'detectable_format:number_highlighted_sections': {'num_highlights': 2}}
            }
        }
    }
    & & \\ \\

    \theutterance \stepcounter{utterance}  
    & & & \multicolumn{2}{p{0.3\linewidth}}{
        \cellcolor[rgb]{0.9,0.9,0.9}{
            \makecell[{{p{\linewidth}}}]{
                \texttt{\tiny{[GM$|$GM]}}
                \texttt{game\_result = LOSE} \\
            }
        }
    }
    & & \\ \\

\end{supertabular}
}

\end{document}
